\documentclass[../../../../main.tex]{subfiles}
\begin{document}

\begin{enumerate}
  \item すべての実数$x$に対して定義された函数
\begin{align*}
  f(x)=\sqrt2a\pi x+\cos\pi x+\sin\pi x \hspace{1\zw} (a\text{は正の定数})
\end{align*}
が極値(極大値または極小値)をもつために，正の定数$a$の満たすべき条件を求めよ．
  \item この条件が満たされているとき，曲線$y=f(x)$の上で，$f(x)$が極大値をとる点は，
すべて1つの直線上に等間隔に並んでいることを示せ．
\end{enumerate}

\end{document}
